\section{Introduction}
\label{sec:intro}

Publicly supported co-creation hubs, enterprise hubs, and related collaborative spaces have become increasingly visible instruments of regional development policy. Local governments often view such hubs as a way to reduce interaction frictions among firms, researchers, entrepreneurs, and other regional actors by creating spaces and programs that facilitate meetings, information exchange, and collaboration. This policy logic is especially attractive in thin regional ecosystems, where the local pool of potentially complementary actors is limited and decentralized coordination may otherwise be weak. Yet the welfare case for hub policy is not obvious in such places. A hub may generate visible early activity while still failing to create enough durable novelty to justify its cost, especially if the same local actors repeatedly interact within a narrow regional pool.

This tension sits naturally within the broader literature on agglomeration and regional coordination. A long tradition in urban and regional economics emphasizes that spatial concentration can raise productivity through matching, learning, and related external economies \citep{DurantonPuga2004,RosenthalStrange2004}. Face-to-face interaction has often been treated as a particularly important channel through which localized economic activity generates value \citep{StorperVenables2004}. Related work on local buzz and global pipelines similarly suggests that geographically proximate interaction can be productive, especially when local contact is complemented by connections to outside knowledge and opportunities \citep{BatheltMalmbergMaskell2004}. More recent evidence from co-working environments also indicates that proximate social interaction can indeed create knowledge spillovers and collaboration gains \citep{RocheOettlCatalini2024}. Taken together, these contributions provide a strong prima facie rationale for why a public hub might help a region.

At the same time, the same literature also implies that more local interaction is not always better. Proximity can support learning, but it can also generate lock-in when interaction becomes repetitive and the same actors repeatedly draw on the same local knowledge base \citep{Boschma2005,TorreRallet2005}. The regional innovation systems literature likewise stresses that local development depends not only on whether interaction occurs, but also on whether regional systems are able to renew themselves through access to external ideas, organizations, and opportunities \citep{CookeUrangaEtxebarria1997,AsheimLawtonSmithOughton2011,TodtlingTrippl2005}. This concern is especially relevant for place-based policy. Local interventions may correct localized coordination problems, but their welfare effects are conditional and context-dependent rather than automatic \citep{KlineMoretti2014ARE,NeumarkSimpson2015}. In the case of public hubs, this raises a simple but underexplored question: when is a publicly supported hub socially efficient in a thin regional ecosystem, and when should it instead be viewed as part of a broader regional policy bundle that also creates external renewal?

This paper develops a simple two-region, two-period theory of public co-creation hubs in thin regional ecosystems. Region $A$ is the target thin region in which a public hub may raise the local matching rate among potential collaborators. Region $B$ is an outside pool from which additional collaborators, ideas, and partnerships may arrive. The model's central distinction is between \emph{local coordination} and \emph{cross-regional renewal}. The hub operates on the first margin by increasing the rate at which potentially valuable local interactions in region $A$ are realized. A complementary \emph{pipeline policy} operates on the second margin by increasing the formation of extra-regional links between region $A$ and outside actors in region $B$.

The model is motivated by a basic dynamic trade-off. In the short run, a hub improves local coordination by lowering interaction frictions and increasing the number of realized local matches. In the longer run, however, that same increase in local interaction intensity may reduce novelty if repeated interaction within a small regional pool becomes redundant. The policy question is therefore not simply whether a hub raises visible local activity. It is whether the hub's short-run coordination gains are large enough to offset fixed cost and the dynamic losses associated with repeated local ties, and whether those losses can be mitigated by policy that expands cross-regional renewal.

The analysis yields three main results. First, a hub is socially efficient as a stand-alone intervention only if its local coordination gain is large enough to offset both its fixed cost and the dynamic redundancy created by repeated local interaction. This produces a simple threshold characterization in terms of ecosystem thickness, matching improvement, fixed cost, and the rate at which local novelty decays. Second, short-run success and social efficiency are not equivalent. The paper identifies a parameter region in which a hub is supported under a private or local short-run evaluation rule because it generates visible first-period activity, yet is socially undesirable once dynamic redundancy is taken into account. Third, the relevant policy object is often not a hub alone, but a \emph{policy bundle} combining local coordination with cross-regional pipeline formation. Even when a stand-alone hub fails a welfare test, the hub may become socially efficient when paired with sufficiently effective outside-linkage policy.

These results speak directly to current regional policy practice. In many settings, hubs are assessed using visible indicators such as participation counts, meetings, events, occupancy rates, or early project volume. The model suggests that such measures are informative about short-run coordination gains, but incomplete as welfare metrics. What matters for social efficiency is not only whether the hub increases local interaction, but also whether interaction remains novel over time and whether the policy design creates adequate channels for outside renewal. In this sense, the paper argues that a public hub in a thin regional ecosystem should not be evaluated as a stand-alone building or venue, but as part of a broader regional coordination architecture.

The paper contributes to three strands of literature. Relative to the agglomeration and interaction literatures, it provides a welfare framework for a policy instrument that is often justified by local buzz and face-to-face contact, but rarely analyzed through the joint lens of coordination gains and dynamic redundancy \citep{DurantonPuga2004,StorperVenables2004,BatheltMalmbergMaskell2004}. Relative to the proximity and regional innovation systems literatures, it connects lock-in and system renewal to an explicit welfare criterion that distinguishes local coordination from cross-regional renewal \citep{Boschma2005,AsheimLawtonSmithOughton2011}. Relative to the place-based policy literature, it shows why the efficient object of evaluation may be a bundled regional policy rather than a stand-alone local project \citep{KlineMoretti2014ARE,NeumarkSimpson2015}.

The model also generates empirically testable implications. It predicts that hub creation should increase initial interaction volume, but not necessarily sustain the share of new ties; that hub-plus-pipeline regions should exhibit stronger outside collaboration and higher medium-run novelty than hub-only regions; and that the same hub policy should be more valuable in regions with greater absorptive capacity. These predictions point toward a more dynamic and network-aware empirical evaluation of hub policy, one that tracks repeated ties, outside linkages, and participant composition rather than relying only on visible activity measures.

The remainder of the paper proceeds as follows. Section~\ref{sec:lit} reviews the related literature and situates the paper's contribution. Section~\ref{sec:model} presents the two-region model. Section~\ref{sec:hub_efficiency} characterizes the conditions under which a hub is socially efficient as a stand-alone policy. Section~\ref{sec:short_run} explains why a place-based public intervention may be needed by distinguishing private short-run evaluation from social welfare evaluation. Section~\ref{sec:refresh} analyzes the optimal pipeline policy and the conditions under which a hub-plus-pipeline bundle is efficient. Section~\ref{sec:extensions} develops extensions on participant composition and digital substitution. Section~\ref{sec:policy} derives empirical predictions and discusses regional policy interpretation. Section~\ref{sec:conclusion} concludes.
